% Source: sources/returns/S6_04_MAPS_MANIFOLDS/code/manifolds_odd_prime_smooth_proof.tex:1–279
\paragraph{Theorem D.2.36 (Proved: the residual five- and seven-primary smooth types).}
Consider the $36$ marked additive-filling patterns on the $30$ initial
rows of D.2.26 with residual $c=5$ or $7$ and no orbifold obstruction.
Use the primitive Mumford cap at every multiplicative position, the
specified minimal original or finite cap at each additive position,
and no additional regular logarithmic transformation. Allow every
positive multiple of the finite monodromy orders and every free
affine twist for which the resulting manifold is simply connected.
For $p=c$, every such manifold is spin and has complete Zhubr data
\[
 H_2=\mathbb Z/p,\qquad b_3=0,\qquad \mu=0,
 \qquad \mathfrak p\in\{0,1\}\quad\text{up to }\operatorname{Aut}(\mathbb Z/p).
 \tag{D.OP1}
\]
Here $4\mathfrak p=\operatorname{PD}(p_1)$, as in D.1.11.
The following entirely local rule determines which value occurs.
There is exactly one zero-weight multiplicative position
whose double curve is fixed by the scalar subgroup $C_p$; it is of
type $I_p$. In its genuine Tate marking put
\[
 (\operatorname{ord}h_1,\operatorname{ord}h_G,
                  \operatorname{ord}R)=(p,r,\rho),\qquad
 w=(ar)^{-1}\pmod p.
\]
Then $\mathfrak p=0$ if and only if $w^2=1$; otherwise its canonical
value is $1$. Thus the $c=5$ patterns give exactly two oriented
diffeomorphism types and the $c=7$ patterns give exactly one.
These types are specified by the full published classifying
invariants (D.OP1), rather than by homology alone.

\begin{proof}
\emph{Marked data and the class which must be extended.}
D.2.26 gives $H_2=H_3=\mathbb Z/p$ and zero middle Betti number
under the stated simple-connectivity condition. Its primitive fibre
class $s$ generates $H_2$. Let $u\in H^2(X;\mathbb F_p)$ evaluate
to one on $s$. In the regular torus lattice the alternating form is
\[
 \Omega=a\,e_1^*e_2^*-b\,\delta^*\gamma^*,\qquad p\mid b,
                                      \quad p\nmid a.
 \tag{D.OP2}
\]
Its values are one on $s$ and zero on the rigidified swept cycles.
The entire finite-cap normal image in (D.HY7) has fibre
coefficients divisible by $b/q_i$, and hence by $p$ in these
patterns. An allowed unchanged additive cap has zero fibre
coefficient, and the ordinary clutch contributes $be$.
Consequently this functional descends modulo $p$ through every
actual cap relation. This is the integral attachment calculation
which identifies its extension with $u$ for all finite orders
and twists; a finite covering transfer is not used.

Remove all additive caps and one ordinary clutch disc, retaining
all multiplicative Mumford caps, and denote the remaining space by
$W$. Its regular part is $V$. The actual multiplicative angular
vanishing columns generate $\ker\gamma^*$ primitively in all
$30$ rows: their Smith factors are $(1,1,1)$. The exact certificate
below uses the full marked columns, including the rank-two cusp.
Van Kampen therefore gives
\[
                  \pi_1(W)=\mathbb Z_\gamma\times F_r
 \tag{D.OP3}
\]
for the free base group with the indicated caps removed. In
particular this group is torsion-free. Each primitive rank-one
Mumford cap and the one-component rank-two cap is generated in
degree two by its boundary. In degree one its sole additional
kernel is the base meridian: the angular lattice is already the
primitive image of the local monodromy minus the identity.
The unfilled ordinary disc makes these peripheral base classes
independent in the graph. Mayer--Vietoris therefore shows that
$H_2(V)\to H_2(W)$ is onto. Together with the free $H_1(W)$ of
(D.OP3), this removes any universal-coefficient ambiguity in
comparing degree-two classes on $W$ from their actual values on $V$.

\emph{A doubled symmetric line.}
The use of $2\Omega$ makes its Appell--Humbert semicharacter
trivial modulo two, irrespective of the parity of $a,b$.
On an auxiliary holomorphic line use the multiplicative factors
\[
\begin{aligned}
 h(z,v,\zeta)&=(h_1z,h_G^a v,z^{-2a}\zeta),\\
 F(z,v,\zeta)&=(h_G^{-b}z,Rv,
                         h_G^{-ab}R^b v^{2b}\zeta),\\
 J(z,v,\zeta)&=(z^{-1},v^{-1},z^{-2a}\zeta).
\end{aligned}                                                    \tag{D.OP4}
\]
Direct substitution gives $hF=Fh$, $J^2=1$, $JhJ=h^{-1}$
and $JFJ=F^{-1}$. The last identity uses precisely the square
of $h_G^{-ab}R^b$. The alternating Chern form is $2\Omega$.
These symmetric line bundles are rigidified at the origin.
Normalized identifications on overlap markings therefore patch:
an automorphism of a rigidified holomorphic line bundle on a
compact torus is the identity. Comparison with the flat real
prequantum family for $2\Omega$ has globally two-torsion Chern
class. Indeed its fibrewise square is canonically trivial after
rigidification, and the unique normalized trivializations patch.
Reduction modulo the odd prime $p$ kills this difference. Thus
half its Chern class modulo $p$ is exactly (D.OP2), including
the zero values on every swept cycle.

We give the required extensions and scalar characters explicitly.
At the rank-two position write the two actual valuation periods as
\[
 H=(n,ar),\quad F=(-br,\rho),\qquad n\rho+ab r^2=1.
\]
An integral support function is
\[
 \psi(x,y)=-a\rho x^2-2ab rxy+bn y^2+ax.
 \tag{D.OP5}
\]
Its differences under $H,F$ are respectively
$-2ax$ and $2by+b\rho-ab r$, the valuations of the
two factors in (D.OP4). The same unimodular fan consequently
extends the line over the cusp. Its vertical slopes
$-2ab rx+bn(2y+1)$ are divisible by $p$.
The scalar $C_p$ character is therefore zero on its fixed
rational double curve $C$. The wall expression, using the two
opposite vertices $(1,0),(-1,1)$ and edge endpoints
$(0,0),(0,1)$, is
\[
 \psi(1,0)+\psi(-1,1)-\psi(0,0)-\psi(0,1)
                         =-2a(\rho-br).
 \tag{D.OP6}
\]
The divisor convention can reverse the sign of this wall
expression; in either convention its half is nonzero modulo
$p$, since $\rho\equiv n^{-1}\pmod p$ and $p\nmid a$.
Thus the class initially normalized by $u(s)=1$ has
$u[C]\ne0$. For the remainder of the proof, rescale $u$
uniquely so that its evaluation on the complex-oriented
curve $C$ is $1$, and scale the auxiliary Chern class by
the same nonzero field element. This fixes the actual
cohomology-generator convention without a support-function
sign ambiguity. In particular $\alpha_C=1$ below.

At a rank-one position write its two periods as
$H=\alpha u,F=\beta u$, with $u=(u_1,u_2)$ primitive and
$\gcd(\alpha,\beta)=1$. Thus $n=\alpha u_1$,
$ar=\alpha u_2$, $-br=\beta u_1$, and $\rho=\beta u_2$.
The identity $a\beta u_1=-b\alpha u_2$ implies
$\alpha\mid au_1$. An explicit integral support is
\[
           \psi_1(j)=-\frac{au_1}{\alpha}\,j(j-\alpha).
 \tag{D.OP5a}
\]
Its differences under $\alpha,\beta$ are exactly
$-2au_1j$ and $2bu_2j+b\rho-ab r$.
The elliptic combination $\alpha F-\beta H$ has zero
exponent in the degenerating coordinate, by the same identity.
These facts extend both specified factors in (D.OP4) over the
primitive fan and retain their actual boundary marking.
At the unique $I_p$ position fixed by $C_p$, the short period
is already the primitive vector $(p,ar)$. Choose
integers $A,B$ with $Ap+B ar=1$. In the monomial coordinates
with first row $(A,B)$ and second row $(-ar,p)$, the scalar
action has normal weight $B=w$ and elliptic weight zero
modulo $p$. Its support function can be taken to be
$-ap\,j(j-1)$. All its slopes are divisible by $p$.
The scalar character of the auxiliary line on this elliptic
double curve $E$ is therefore also zero.

These are equivariant extensions. On the covering coordinates
the scalar subgroup lifts by $(z,v,\zeta)\mapsto
(z,\lambda v,\zeta)$ for $\lambda^p=1$; it commutes with
(D.OP4) because $p\mid b$. In the prequantum comparison the
same lift is preserved by every linear marking: each fixes
$\delta$ and $2\Omega$, and the semicharacter is trivial.
The symmetric two-torsion comparison has a unique compatible
odd-order lift. Holomorphic units on the local base discs
have logarithms and contribute no scalar character or wall
degree. Thus the above local lifts patch on $W$.

\emph{The fixed set and the global equivariant comparison.}
The primitive rank-one period at a multiplicative cap is the
column $(u_1,u_2)$ obtained by taking the Hermite basis of its
two period columns. Scalar isotropy along its double curve
has order $|u_1|$. The retained exact data show that precisely
one such order is divisible by $p$, and its type is $I_p$.
The rank-two fan has just the vertical rational fixed curve
$C$; its other primitive edge directions are not parallel to
$\delta$ modulo $p$. There are no further fixed strata.
An allowed original additive cap is a product with the free
scalar circle. At a finite cap, primitivity of the actual gamma
numerator in a simply connected member excludes scalar fixed
points: equality of a scalar translate with the $j$th affine
power would force $m\mid j$ from the gamma coordinate, and
then require a nontrivial scalar translation to fix a torus
point. This also covers orders divisible by $p$.
Consequently
\[
                         X^{C_p}=C\sqcup E.
 \tag{D.OP7}
\]

The class $u$ admits a Borel equivariant lift. The obstruction
with $H^1(X;\mathbb F_p)$ vanishes by simple connectivity;
the remaining differential into $H^3(BC_p;\mathbb F_p)$
vanishes because any point of $C$ splits the base map.
The whole scalar circle fixes that point and acts trivially
on the based fundamental group of $W$. Hence
\[
 \pi_1(W_{C_p})=(\mathbb Z\times F_r)\times C_p.
\]
Every fixed point determines its same, unique order-$p$
subgroup, by torsion-freeness of the first factor.
The difference between the chosen lift of $u$ and the correspondingly scaled half of
the equivariant Chern class of the auxiliary line has zero
ordinary restriction on $W$. The universal-cover low-degree
sequence places this difference in group cohomology.
Its restrictions to the fixed-point copies of $BC_p$ are
therefore equal. Adjust the lift once by a base class.
Writing $t\in H^2(BC_p;\mathbb F_p)$ and
$\eta\in H^1(BC_p;\mathbb F_p)$, we obtain
\[
 \widetilde u|_C=\alpha_C[C]^*,\qquad
 \widetilde u|_E=\alpha_E[E]^*+\eta\zeta,
                  \quad\zeta\in H^1(E;\mathbb F_p).
 \tag{D.OP8}
\]
In particular both restrictions have zero coefficient of $t$.

\emph{All the classifying invariants.}
Since $H^2(X;\mathbb Z)=0$, the complex first Chern class
vanishes and $X$ is spin. The two complex normal weights
at $C$ are $1,-1$, and those at $E$ are $w,-w$.
One can also check the sums from equivariant $c_1$:
it comes from $H^2(BC_p;\mathbb Z)$, and is zero at $C$,
hence zero everywhere. The equivariant Euler classes
have leading terms $-t^2$ and $-w^2t^2$.
Inverting $t$ in the equivariant Thom isomorphism gives
localization on (D.OP7). Applying it to $u$, whose integral
has negative degree and is zero, gives
\[
                      \alpha_C+\alpha_E/w^2=0.
 \tag{D.OP9}
\]
No mixed term in (D.OP8) contributes: it has ordinary
degree one and its product with an ordinary degree-two
Euler correction has degree three on a curve.

For each normal bundle with roots $wt+x_1,-wt+x_2$,
the square of their sum is zero on the fixed curve.
The tangent line's square is zero as well. Therefore its
equivariant Pontryagin class divided by the normal Euler
class is exactly $-2$. It follows that
\[
 \langle u\,p_1(X),[X]\rangle
   =-2(\alpha_C+\alpha_E)
   =-2(1-w^2).
 \tag{D.OP10}
\]
Moreover the cube of each restriction in (D.OP8) is zero:
ordinary degree-two classes square to zero on a curve,
and $\eta^2=0$ at an odd prime. Thus
\[
                        \langle u^3,[X]\rangle=0.
 \tag{D.OP11}
\]
Lemma D.1.10 with $r=0,d=p$ identifies
$H_6(K(\mathbb Z/p,2);\mathbb Z)=\mathbb Z/p$
and proves that (D.OP11) detects its entire fundamental
class invariant. Consequently $\mu=0$.
Multiplication by four is invertible on $H_2$, so
(D.OP10) computes the entire refined Pontryagin class.
Because $\alpha_C\ne0$, it is zero exactly when $w^2=1$.
Every nonzero element of $\mathbb Z/p$ is carried to one
by an automorphism; these automorphisms leave $\mu=0$.
Zhubr's Theorem~1, with the definitions on pp.~794--795
\cite{manZhubr75}, now gives precisely (D.OP1) and the
claimed oriented diffeomorphism classification. Existence
comes from the actual constructed members, so no separate
realization assumption is needed.

The local weights in the retained table include both
$w^2=1$ and $w^2\ne1$ for $p=5$, and only $w^2\ne1$
for $p=7$. The code checks all $30$ actual marked rows;
the $36$ additive patterns and all allowed higher orders
use the same local data. This proves the final counts.
\end{proof}
